\subsection{Faithfully flat rings}

PROPOSITION 8. Let A, B be two rings and \(\rho\) a homomorphism from A to B. Suppose that there exists a right B-module E such that \(\rho_{*}(E)\) is a faithfully flat A-module. Then:

\begin{enumerate}
\def\labelenumi{(\roman{enumi})}
\item
  For every left A-module F, the canonical homomorphism \(j: \mathbf{F} \to \mathbf{F}_{(\mathbf{B})} = \mathbf{B} \otimes_{\mathbf{A}} \mathbf{F}\) (such that \(j(x) = 1 \otimes x\) for all \(x \in \mathbf{F}\) ) is injective.
\item
  For every left ideal \(\mathfrak{a}\) of \(\mathbf{A},\overline{\rho}^{-1}(\mathbf{Ba}) = \mathfrak{a}\) .
\item
  The homomorphism \(\rho\) is injective.
\item
  For every maximal left ideal \(\mathfrak{m}\) of A, there exists a maximal left ideal \(\mathfrak{n}\) of B such that \(\overline{\rho}^{-1}(\mathfrak{n}) = \mathfrak{m}\) .
\end{enumerate}

We prove (i). We know (Algebra, Chapter II, § 5, no. 2, Corollary to Proposition 5) that for every right B-module M, the canonical A-homomorphism \(i: \mathbf{M} \to \rho_*(\mathbf{M}) \otimes_{\mathbf{A}} \mathbf{B} = \rho^*(\rho_*(\mathbf{M}))\) defined by \(i(y) = y \otimes 1\) is injective and that the A-module \(i(\mathbf{M})\) is a direct factor of \(\rho_*(\mathbf{M}) \otimes_{\mathbf{A}} \mathbf{B}\) . Hence, for every left A-module F,

\[
i \otimes 1 _ {F}: \rho_ {*} (M) \otimes_ {A} F \rightarrow \rho_ {*} (M) \otimes_ {A} B \otimes_ {A} F
\]

is injective (§ 2, no. 1, Lemma 2). Taking \(\mathbf{M} = \mathbf{E}\) , it follows (since \(i \otimes 1_{\mathbf{F}} = 1_{\mathbf{M}} \otimes j\) ) that \(j\) is injective (no. 1, Proposition 2).

Assertion (ii) follows from (i) by taking \(\mathbf{F} = \mathbf{A}_s / \mathfrak{a}\) and (iii) from (ii) by taking \(\mathfrak{a} = \{0\}\) .

Finally, if \(\mathfrak{m}\) is a maximal left ideal of \(\mathbf{A}\) , then \(\bar{\rho}^{-1}(\mathbf{B}\mathfrak{m}) = \mathfrak{m}\) by (ii), and consequently \(\mathbf{B}\mathfrak{m} \neq \mathbf{B}\) . Then there exists a maximal left ideal \(\mathfrak{n}\) of \(\mathbf{B}\) containing \(\mathbf{B}\mathfrak{m}\) (Algebra, Chapter I, § 8, no. 7, Theorem 2); then \(\mathfrak{m} \subset \bar{\rho}^{-1}(\mathfrak{n})\) and as \(\rho(1) \notin \mathfrak{n}\) , \(1 \notin \bar{\rho}^{-1}(\mathfrak{n})\) . Consequently \(\bar{\rho}^{-1}(\mathfrak{n}) = \mathfrak{m}\) .

If A and B satisfy the conditions of Proposition 8, A is usually identified with a subring of B by means of \(\rho\) .

COROLLARY. Under the hypotheses of Proposition 8, if B is left Noetherian (resp. Artinian), so is A.

If \((\mathfrak{a}_n)\) is a non-stationary increasing (resp. decreasing) sequence of left ideals of A, the sequence \((\mathrm{Ba}_n)\) of ideals of B is non-stationary and increasing (resp. decreasing) since \(\overline{\rho}^{-1}(\mathrm{Ba}_n) = \mathfrak{a}_n\) , contrary to the hypothesis.

*Remark (1). If A and B are commutative, we shall see in Chapter II, § 2, no. 5, Corollary 4 to Proposition 11, that the hypothesis of Proposition 8 implies that

for every \(\text{prime}\) ideal \(\mathfrak{p}\) of \(\mathbf{A}\) , there exists a \(\text{prime}\) ideal \(\mathfrak{q}\) of \(\mathbf{B}\) such that \(\overline{\rho}^{-1}(\mathfrak{q}) = \mathfrak{p}\) (where \(\mathfrak{p} = \mathbf{A} \cap \mathfrak{q}\) when \(\mathbf{A}\) is identified with a subring of \(\mathbf{B}\) ).*

An important application of Proposition 8 is when B is itself a faithfully flat A-module. But in this case we have the following more precise proposition:

PROPOSITION 9. Let A, B be two rings and \(\rho\) a homomorphism of A to B. The following five properties are equivalent:

\begin{enumerate}
\def\labelenumi{(\alph{enumi})}
\item
  The right A-module B is faithfully flat.
\item
  The homomorphism \(\rho\) is injective and the right A-module \(\mathbf{B} / \rho (\mathbf{A})\) is flat.
\item
  The right A-module B is flat and, for every left A-module F, the canonical homomorphism \(x \mapsto 1 \otimes x\) of F to B \(\otimes_{A} F\) is injective.
\item
  The right A-module B is flat and, for every left ideal a of A, \(\overline{\rho}^{-1}(\mathbf{B}\mathfrak{a}) = \mathfrak{a}\) .
\item
  The right A-module B is flat and, for every maximal left ideal m of A, there
\end{enumerate}

exists a maximal left ideal n of B such that \(\bar{\rho}^{-1}(\mathfrak{n}) = \mathfrak{m}\) .

By Proposition 8, (a) implies each of the properties (c), (d), (e). On the other hand, if (e) holds, then \(\mathbf{Bm} \neq \mathbf{B}\) for every maximal left ideal \(m\) of A (since there exists a maximal left ideal \(n\) of B such that \(\mathbf{Bm} \subset n\) ), and B is a faithfully flat A-module by criterion (d) of Proposition 1 of no. 1; hence (e) implies (a).

We shall now prove that \((c) \Rightarrow (d) \Rightarrow (b) \Rightarrow (a)\) , which will complete the proof. In the first place, (c) implies (d) by taking \(F = A_{s}/a\) in (c). If (d) holds, by taking \(a = \{0\}\) it follows that \(\rho\) is injective; (d) and §2, no. 6, Corollary to Proposition 7 imply that \(B/\rho(A)\) is a flat right A-module, that is (d) implies (b). Finally, if (b) holds, Proposition 3 of no. 1 applied to the exact sequence

\[
0 \longrightarrow \mathrm{A} _ {d} \stackrel {\rho} {\longrightarrow} \mathrm{B} \longrightarrow \mathrm{B} / \rho (\mathrm{A}) \longrightarrow 0
\]

shows that B is a faithfully flat right A-module, since \(A_{d}\) is faithfully flat.

*Remark (2). If A and B are commutative, we shall see in Chapter II, § 2 no. 5, Corollary 4 to Proposition 11 that the conditions of Proposition 9 are equivalent to the following:

\begin{enumerate}
\def\labelenumi{(\alph{enumi})}
\setcounter{enumi}{5}
\tightlist
\item
  B is a flat A-module and, for every prime ideal p of A, there exists an ideal q of B such that \(\overline{\rho}^{-1}(\mathfrak{q}) = \mathfrak{p}\) .
\end{enumerate}

Under the conditions of Proposition 9, let us identify A with a subring of B by means of \(\rho\) . The relation \(\bar{\rho}^{-1}(\mathbf{Ba}) = a\) then reads \(\mathbf{A} \cap \mathbf{Ba} = a\) . On the other hand if F is a left A-module, F is identified with its image in \(\mathbf{B} \otimes_{\mathbf{A}} \mathbf{F}\) under the canonical mapping \(x \mapsto 1 \otimes x\) ; if X is an additive subgroup of F, then we denote by BX the left sub-B-module of \(\mathbf{B} \otimes_{\mathbf{A}} \mathbf{F}\) generated by X. With this notation, we have:

PROPOSITION 10. Let B be a ring and A a subring of B such that B is a faithfully flat right A-module. Let F be a left A-module, F', F'' two submodules of F. Then:

\begin{enumerate}
\def\labelenumi{(\roman{enumi})}
\item
  The canonical mapping \(\mathbf{B} \otimes_{\mathbf{A}} \mathbf{F}' \to \mathbf{B} \otimes_{\mathbf{A}} \mathbf{F}\) is an isomorphism of \(\mathbf{B} \otimes_{\mathbf{A}} \mathbf{F}'\) onto \(\mathbf{BF}'\) .
\item
  \(\mathbf{F} \cap \mathbf{BF}' = \mathbf{F}'\) .
\item
  \(\mathrm{B}(\mathbf{F}' + \mathbf{F}'') = \mathbf{BF}' + \mathbf{BF}''\) .
\item
  \(\mathbf{B}(\mathbf{F}'\cap \mathbf{F}'')=\mathbf{B}\mathbf{F}'\cap \mathbf{B}\mathbf{F}''\) .
\end{enumerate}

As B is a flat right A-module, the canonical mapping

\[
\mathbf {B} \otimes_ {\mathbf {A}} \mathbf {F} ^ {\prime} \rightarrow \mathbf {B} \otimes_ {\mathbf {A}} \mathbf {F}
\]

is injective; taking account of the identifications made, its image is BF', which proves (i). Assertion (ii) follows from § 2, no. 6, Proposition 7, applied with E = B, E' = A, and using the formulae \(A \otimes_{A} F = F\) and \(A \otimes_{A} F' = F'\). Assertion (iii) is trivial and (iv) follows from § 2, no. 6, Proposition 6.

